% Job posting: https://ca.linkedin.com/jobs/view/computer-vision-software-developer-professional-ii-d%C3%A9veloppeur-ou-d%C3%A9veloppeuse-de-logiciels-en-vision-par-ordinateur-ii-at-zebra-technologies-4463076728
\documentclass[11pt]{article}
\usepackage[left=1in,top=1in,right=1in,bottom=1in]{geometry}
\usepackage[hidelinks]{hyperref}
\usepackage{setspace}
\usepackage[T1]{fontenc}
\usepackage{newtxtext,newtxmath}
\ifdefined\pdfgentounicode
  \input{glyphtounicode}
  \pdfgentounicode=1
\fi
\setstretch{1.1}
\pagenumbering{gobble}
\begin{document}
\pagestyle{empty}

\noindent
Nishal Stanislaus Silva, PhD \\
Toronto, ON \\
nishal.silva@hotmail.com \,|\, +1 613 200 2030 \\
\href{https://nishal.xyz}{nishal.xyz} \,|\, \href{https://www.linkedin.com/in/nishal-silva}{linkedin.com/in/nishal-silva} \,|\, \href{https://github.com/ns2max}{github.com/ns2max}

\vspace{1em}
\noindent Dear Hiring Manager,

\vspace{1em}
\noindent For 5.5 years at MAS Holdings, South Asia's largest apparel manufacturer, I built and deployed computer-vision systems directly onto production floors. I automated multilingual care-label quality control using OpenCV template and feature matching with an explainable defect overlay, iterating the system from a benchtop scanner to an industrial camera to a full conveyor line with PLC-driven reject logic, and cutting inspection time by 99.5\% while keeping a human-in-the-loop path for borderline cases. On a separate line, I built a loom-side fabric structural-defect detection system using a microscopic camera array that let a pilot deployment cut operator headcount by 50\%. Zebra builds the hardware and software that make exactly this kind of production-floor inspection possible at scale, and I would like to bring that same track record to your Computer Vision Software Developer team in Montreal.

\vspace{1em}
\noindent What I bring is the full path from prototype to deployed production, not just a model in a notebook. Beyond the QC work, I designed the camera-captured-garment-geometry-to-warp alignment algorithm for Promptly, MAS's on-demand direct-to-garment printing system now running commercially in the US, Mexico, France and Sri Lanka, and personally built its RIP integration. I write production C++17 and Python, have introduced CI/CD and code review into an engineering group, and have deployed inference on embedded and ARM-class hardware under real latency and compute budgets. My PhD in Information Engineering and Computer Science from the University of Trento added a further layer of rigor: designing frozen-protocol benchmarks and ablations to quantify accuracy, latency and CPU trade-offs before a system ships, a habit I have carried into every project since.

\vspace{1em}
\noindent I have not worked directly with barcode or RFID scanning hardware, but the core discipline this role calls for, taking a camera-to-decision pipeline from prototype to a reliable, embedded, production-grade system on a factory floor, is exactly what I have spent 5.5 years doing, and I would expect that gap to close quickly.

\vspace{1em}
\noindent I am a Canadian Permanent Resident and can work in Canada without sponsorship. I am based in Toronto and open to Montreal, where I spent a visiting research term at McGill's IDMIL/CIRMMT lab and already have ties to the local research and engineering community. I am available immediately and would welcome the opportunity to discuss how my industrial computer-vision background can contribute to Zebra's team.

\vspace{1em}
\noindent Sincerely, \\
Nishal Stanislaus Silva

\end{document}
